\documentclass[11pt]{article}
\usepackage{amsmath}
\usepackage{enumitem}
\usepackage{geometry}
\usepackage{hyperref}
\geometry{margin=1in}

\title{EC 441 -- Week Review \\ Multiple Access Protocols and Ethernet}
\date{}
\begin{document}
\maketitle

\section*{Overview}
This week covered how nodes share a broadcast channel (MAC protocols) and how Ethernet implements link-layer delivery using frames, MAC addresses, ARP, and switching.

\hrule
\vspace{0.5cm}

%%%%%%%%%%%%%%%%%%%%%%%%%%%%%%%%%%%%%%%%%%%%%%%%%%%%%%%%%%
\section{Lecture 7: Multiple Access Protocols}

\subsection{The Multiple Access Problem}
On a shared broadcast medium (e.g., classic Ethernet, WiFi, satellite), multiple nodes may transmit simultaneously, causing collisions.  
A MAC protocol coordinates who transmits and when.

\subsection{Channel Partitioning Protocols}

\subsubsection*{TDMA (Time Division Multiple Access)}
\begin{itemize}
    \item Time divided into slots.
    \item Each node gets one fixed slot per frame.
    \item No collisions.
    \item Inefficient when nodes are idle.
\end{itemize}

\subsubsection*{FDMA (Frequency Division Multiple Access)}
\begin{itemize}
    \item Frequency band divided into sub-bands.
    \item Each node permanently assigned one band.
    \item Suffers from idle bandwidth waste.
\end{itemize}

\textbf{Key limitation:} Partitioning wastes capacity under bursty traffic.

\subsection{Taking Turns}

\subsubsection*{Polling}
\begin{itemize}
    \item Central coordinator polls nodes.
    \item No collisions.
    \item Single point of failure.
\end{itemize}

\subsubsection*{Token Passing}
\begin{itemize}
    \item Token circulates in a logical ring.
    \item Only token holder may transmit.
    \item Predictable latency.
\end{itemize}

\subsection{Random Access Protocols}

\subsubsection*{Pure ALOHA}
\begin{itemize}
    \item Transmit immediately.
    \item Retransmit after random delay if collision.
    \item Collision window: $2T$.
\end{itemize}

Throughput:
\[
S = G e^{-2G}
\]
Maximum:
\[
S_{max} = \frac{1}{2e} \approx 0.18
\]

\subsubsection*{Slotted ALOHA}
\begin{itemize}
    \item Transmissions start at slot boundaries.
    \item Collision window reduced to $T$.
\end{itemize}

Throughput:
\[
S = G e^{-G}
\]
Maximum:
\[
S_{max} = \frac{1}{e} \approx 0.37
\]

\subsection{CSMA (Carrier Sense Multiple Access)}
\begin{itemize}
    \item Listen before transmitting.
    \item Collisions only occur within propagation delay $\tau$.
    \item Much higher efficiency than ALOHA.
\end{itemize}

\subsection{CSMA/CD (Collision Detection)}
\begin{enumerate}
    \item Sense channel.
    \item Transmit if idle.
    \item Detect collision while transmitting.
    \item Send jam signal.
    \item Binary exponential backoff.
\end{enumerate}

Binary exponential backoff:
\[
\text{Wait from } \{0,1,\dots,2^{\min(k,10)}-1\}
\]

\subsection{Minimum Frame Size Constraint}

To detect collisions:
\[
T \ge 2\tau
\]

Since $T = \frac{L}{R}$:
\[
L \ge 2\tau R
\]

Classic Ethernet:
\[
L_{min} = 64 \text{ bytes}
\]

Modern switched Ethernet is full-duplex, so CSMA/CD is not used.

\subsection{Wireless Limitation}
Collision detection fails in wireless due to:
\begin{itemize}
    \item Self-interference.
    \item Hidden terminal problem.
\end{itemize}

WiFi uses CSMA/CA (collision avoidance) instead.

\vspace{0.5cm}
\hrule
\vspace{0.5cm}

%%%%%%%%%%%%%%%%%%%%%%%%%%%%%%%%%%%%%%%%%%%%%%%%%%%%%%%%%%
\section{Lecture 8: Ethernet, Addressing, Switching, ARP}

\subsection{Ethernet Frame Structure}

\begin{center}
Preamble | Dst MAC | Src MAC | Type | Payload | FCS
\end{center}

\begin{itemize}
    \item Preamble (7B) + SFD (1B)
    \item Destination MAC (6B)
    \item Source MAC (6B)
    \item Type (2B)
    \item Payload (46--1500B)
    \item FCS (4B, CRC-32)
\end{itemize}

\textbf{Minimum frame size:} 64 bytes  
\textbf{Maximum frame size:} 1518 bytes

\subsection{MAC Addresses}

48-bit identifier written in hex:
\[
\texttt{A4:C3:F0:85:AC:2D}
\]

\begin{itemize}
    \item First 3 bytes: OUI (manufacturer).
    \item Last 3 bytes: device-specific.
    \item Broadcast: FF:FF:FF:FF:FF:FF
\end{itemize}

MAC addresses are \textbf{flat}.  
IP addresses are \textbf{hierarchical}.

\subsection{ARP (Address Resolution Protocol)}

Purpose: Resolve IP address to MAC address within a subnet.

\textbf{ARP Request:}
\begin{itemize}
    \item Broadcast frame.
    \item "Who has IP X?"
\end{itemize}

\textbf{ARP Reply:}
\begin{itemize}
    \item Unicast back to requester.
    \item Contains MAC address.
\end{itemize}

Entries stored in ARP cache with TTL.

\subsection{Subnets}

CIDR notation:
\[
192.168.1.0/24
\]

\begin{itemize}
    \item First 24 bits: network prefix.
    \item Remaining bits: host portion.
\end{itemize}

ARP works only within a subnet.  
Cross-subnet traffic requires a router.

\subsection{Switching and Self-Learning}

Switch builds a forwarding table:

\begin{center}
MAC Address $\rightarrow$ Port
\end{center}

Algorithm:
\begin{enumerate}
    \item Learn source MAC on incoming port.
    \item If destination known → forward to specific port.
    \item If unknown → flood.
    \item Broadcast → flood.
\end{enumerate}

Entries expire after inactivity.

\subsection{Hub vs Switch}

\textbf{Hub:}
\begin{itemize}
    \item Repeats to all ports.
    \item Single collision domain.
\end{itemize}

\textbf{Switch:}
\begin{itemize}
    \item Forwards selectively.
    \item Each port is its own collision domain.
    \item Full-duplex.
\end{itemize}

\subsection{Ethernet Evolution}

10 Mb/s → 100 Mb/s → 1 Gb/s → 10/100/400 Gb/s

Line codes evolved:
\begin{itemize}
    \item Manchester (10 Mb/s)
    \item 4B5B + MLT-3 (100 Mb/s)
    \item PAM-5 (1 Gb/s)
    \item PAM-4 + DSP (10+ Gb/s)
\end{itemize}

\section*{Big Picture Connections}

\begin{itemize}
    \item Lecture 7 explains \textbf{why 64-byte minimum frame exists}.
    \item Lecture 8 shows how switching removed collisions but kept frame format.
    \item ARP connects network layer (IP) to link layer (MAC).
    \item Modern Ethernet: switched, full-duplex, high-speed, backward-compatible.
\end{itemize}

\end{document}